\section{Conclusion}
Reference-preserving speculative execution requires a qualified transaction, not merely correct acceptance logic at one decision. The studied RTX implementation constrains attention query layout during target prefill and continuation while retaining other batched computation. Natural crosses and a four-switch qualification study support that finite repair scope; complete responses and width failures define its boundary. Paid ordinary and matched comparisons then separate the deployment gain from the speculative increment. The resulting decisions are mixed: the frozen RTX anchor helps on aggregate, a wider qualified draft loses, the B570 repair remains slower than strong baselines, and M4 and 358H restrict the case for added control and draft-only optimization. The study supplies concrete implementation evidence and reproducible cost comparisons within one model family, while preserving failures that prevent broader claims.
